% DRAFT, for a later decision (brief of 25 Nov 2026, item 4). Not part of paper P5. Not posted.
% Written 25 Nov 2026 by a0805369bae2ebddc (Claude Opus 5.5 at effort high).
% Status of each item and who checked what: see the comments before each section.
% An e-mail on items 1 and 2 went to K. D. Mulmuley on 24 Nov 2026 (reports_20261124/EMAILS_SENT_20261124.md, section 6);
% wait for his reply before this text is final. Item 3 (N-012) was not in that e-mail.
\documentclass[11pt,a4paper]{article}
\usepackage[T1]{fontenc}
\usepackage[utf8]{inputenc}
\usepackage{lmodern}
\usepackage{amsmath,amssymb,amsthm}
\usepackage[margin=1.15in]{geometry}
\usepackage[hidelinks]{hyperref}
\newcommand{\Q}{\mathbb{Q}}
\newtheorem{proposition}{Proposition}
\theoremstyle{definition}
\newtheorem{remark}{Remark}

\title{Three corrections to ``Geometric complexity theory VIII''
(draft)\thanks{AI-produced; see the end for who found and checked what. No
human has checked this text. To the extent that any rights exist in it, they
are dedicated to the public domain under CC0~1.0.}}
\author{}
\date{25 November 2026 (draft)}

\begin{document}
\maketitle

\noindent We refer to K.~D.~Mulmuley, \emph{Geometric complexity theory VIII:
On canonical bases for the nonstandard quantum groups}, arXiv:0709.0751v2
(GCT~VIII below). Locators are those of that version; the writer of this draft
has not opened it (see the end).

% Item 1. Found 21 Nov 2026 by Claude Opus 5.5 at effort high (computation of all 1000 structure constants of
% E^opp(3), both sign readings); closed-form proof by GPT-6 Astra (ChatGPT Pro mode), 23-24 Nov; product identities
% verified by Claude Fable 5.2 at effort max on 24 Nov (left regular representation, exact arithmetic), which also
% compared the statements with the source. Scalar identities below rechecked by this writer (checks/errata_viii_scalar_check.py).
\section{The unimodality clause fails at rank three}

In the section ``Kronecker problem: $n=2$, $r=3$'', the subsection
``Positivity'' states that every multiplicative structure constant of the
basis $E^{\mathrm{opp}}(3)$ has the form $(-1)^a(q^{1/2}-q^{-1/2})^bf$ with $f$
positive and unimodal. One constant printed in the figure of that subsection
is
\[
 c(\gamma^2_{21};\gamma^2_{21};\gamma^2_{21})
 =\frac{(1+q^2)^2(q^2+q+1)(q-1)^2}{q^4}
 =(q^{1/2}-q^{-1/2})^2f(q),
\]
\[
 f(q)=(q+q^{-1})^2(q+1+q^{-1})=q^3+q^2+3q+2+3q^{-1}+q^{-2}+q^{-3}.
\]
The coefficients $1,1,3,2,3,1,1$ are not unimodal ($f_{-1}=3>f_0=2$), in the
sense defined after the first Positivity conjecture (Section~2, ``Positivity in
the Kronecker case''). The exponent and sign cannot be chosen otherwise: the
constant vanishes to order exactly two at $q=1$. Six structure constants of
$E^{\mathrm{opp}}(3)$ have this $f$. All constants are of the form
$\pm(q^{1/2}-q^{-1/2})^af$ with $f$ positive and invariant under $q\mapsto
q^{-1}$; only unimodality fails.

\paragraph{A short derivation.}
Put $D=(q+1)^2/q$, $L=(q+q^{-1})^2$ and $X=P_{21}$. With $b_2=-D$ (see
Section~\ref{sec:b2}), $\gamma^2_1=P_1(X-D)$, $\gamma^2_{21}=X(X-D)$,
$\Sigma=P_1(X-D)(X-L)$, and $\Sigma X=D^2\Sigma$. From
$X(X-D)-L(L-D)=(X-L)(X+L-D)$ one gets
\[
 \gamma^2_1\gamma^2_{21}=L(L-D)\gamma^2_1+(D^2+L-D)\Sigma,\qquad
 L(L-D)=(D-2)^2(D-1)(D-4),
\]
with $D-4=(q^{1/2}-q^{-1/2})^2$ and $D-1=q+1+q^{-1}$. The same coefficient
$L(L-D)$ is the constant $c(\gamma^2_{21};\gamma^2_{21};\gamma^2_{21})$ above.

% Item 2. Same finders and checkers as item 1.
\section{The sign of \texorpdfstring{$b_2$}{b2}}\label{sec:b2}

The display of $E^{\mathrm{opp}}(3)$ uses $b_1$ in the rows of $\gamma^i$ for
both $i=1$ and $i=2$, and gives $b_2=(q+1)^2/q$. The constants in the figure of
the ``Positivity'' subsection are those of
$\gamma^2_x=-((q+1)^2/q)P_x+P_{x'}$, that is, with $b_2$ negative and used in
the second family. With $b_2$ positive, as printed, the ten elements again form
a basis, and the coefficient of $\gamma^2_1$ in $\gamma^2_1\gamma^2_{21}$ is
$(q+q^{-1})^2(q^2+q+4+q^{-1}+q^{-2})$, with coefficients
$1,1,6,3,10,3,6,1,1$, again not unimodal.

% Item 3. GCT_POS run (GPT-6 Astra, ChatGPT Pro mode), result N-012, file
% research_returns/GCT_POS_final_20261122T0219Z_unzipped/branches/nonstandard/viii_unbounded_scaling.tex.
% The run's own reviews read the proof and compared it with the primary PDF (CHECK_EVENTS rows 26, 146, 191, 192).
% Board row H23 records "N-012 proof checked"; the checker is not named there. This writer read the run's text on
% 25 Nov and checked the rank-one example below by hand; it did not open GCT VIII. The rank-four extension rests on
% the run's result N-010 (rational GT lines), which this writer has not read.
\section{Conjecture 3.2 as displayed has no minimum}

GCT~VIII, Subsection 3.3, Conjecture 3.2 (printed p.~22), varies the rational
exponent $a$ in $c_{a,T}=q^ac'_T$ and asks for an element of minimal ambient
degree in the fibre
\[
 \psi_{a,T}^{-1}(b_{a,T})\subset R^{a,\le T}\cap\bar R^{a,\le T}\cap
 \mathcal B^{\le T}_{\Q},
\]
where $b_{a,T}$ is the class of $c_{a,T}$ modulo $q^{1/2}R^{a,\le T}$ and the
degree of a coefficient is its pole order at $q=\infty$ (Subsection 3.3,
equation (13), p.~21). The displayed domain is not intersected with the
bar-fixed elements.

\begin{proposition}
Under these definitions the minimum need not exist. Already for the
one-dimensional algebra $\mathcal B=\Q(q^{1/2})$ with $c'_T=1$ and Laurent form
$\Q[q^{1/2},q^{-1/2}]$, the elements $y_n=q^{-n}$, $n\ge0$, are admissible for
$a_n=-n$, and their degrees $-n$ are unbounded below.
\end{proposition}

\begin{proof}
Let $A_0$, $A_\infty$ be the rings of functions regular at $0$ and $\infty$.
For $a_n=-n$, $R^{a_n}=q^{-n}A_0$ and $\bar R^{a_n}=q^nA_\infty$. The element
$y_n=q^{-n}$ lies in the Laurent form and in $R^{a_n}$, and its quotient by the
generator $q^{-n}$ of $R^{a_n}$ is $1$, so it maps to $b_{a_n,T}$. Its quotient
by the generator $q^n$ of $\bar R^{a_n}$ is $q^{-2n}$, regular at infinity, so
$y_n\in\bar R^{a_n}$. Its only coefficient has degree $-n$. A nonzero algebraic
function has a finite pole order, so no candidate has degree at most $-n$ for
all $n$.
\end{proof}

\begin{remark}
The run that found this (GCT\_POS, result N-012) also shows that the same family
$q^{-n}t$, with $t$ a bar-fixed element on a minimal Gelfand--Tsetlin line, is
admissible in the rank-four algebra for every label order and every initial
normalization, using its rational Gelfand--Tsetlin construction (N-010). One
repair is to require $y=\bar y$ in the displayed domain, as in GCT~VIII,
Section~4, which starts from a bar-invariant element; the counterexample does
not touch that internal procedure.
\end{remark}

\section*{Who found and checked what}
Items 1 and 2 were found on 21 November 2026 by Claude Opus~5.5 at effort high,
which computed all 1000 structure constants of $E^{\mathrm{opp}}(3)$ in
$H_3\otimes H_3$ in both readings of the sign; they were proved again in closed
form by GPT-6 Astra (ChatGPT Pro mode) on 23--24 November; Claude Fable~5.2 at
effort max verified both product identities by a third method on 24 November and
compared the statements with the source. Item 3 is result N-012 of a run of
GPT-6 Astra (ChatGPT Pro mode), 21 November 2026. No human has checked any of
this.

\end{document}
